\begin{lemma}\label{rl-000U}\label{lem:euler-class-well-defined}
  Let \(T\) be a split diagonalizable group and let \(\mathcal X\) be an algebraic stack with trivial \(T\)-action. Let
  \begin{align*}
    E^\bullet = [E^0 \to E^1], \hspace{1em} F^\bullet = [F^0 \to F^1]
  \end{align*}
  be a two-term complexes of finite-rank moving \(T\)-equivariant locally free sheaves. If \(E^\bullet\) and \(F^\bullet\) are isomorphic in the equivariant derived category, then
  \begin{align*}
    e_T(E^1)e_T(F^0) = e_T(E^0)e_T(F^1)
  \end{align*}
  in \(A^*_T(\mathcal X)\). In particular, \(e_T(E^\bullet) := e_T(E^1)e_T(E_0)^{-1}\) and \(e_T(F^\bullet)\) agree in \(A_*^T(\mathcal X)_{\loc}\)
\end{lemma}

\begin{proof}
  Fix a complex \(G_\bullet = [G_0 \xrightarrow{d} G_1]\) of quasicoherent sheaves and morphisms \(\phi:E_\bullet \to G_\bullet\) and \(\psi:F_\bullet\to G_\bullet\) exhibiting the isomorphism. Given the data of a locally free sheaf \(A_1\) and a morphism \(\chi:A_1\to G_1\) such that the induced map \(\overline \chi: A_1 \to H^1(G_\bullet)\) is surjective, we can find a two-term complex \(A^\chi_\bullet\) of locally free sheaves which is quasi-isomorphic to \(G_\bullet\). We do this because while we can't take the Euler class of \(G_\bullet\), we can take the Euler class of \(A_\bullet^\chi\).

  First, consider the map \(\chi - d: A_1\oplus G_0 \to G_1\) and set \(A_0 = \ker (\chi - d)\). Then define \(A^\chi_\bullet = [A_0 \xrightarrow{d_A} A_1]\), where \(d_A\) is the composition \(A_0 \hookrightarrow A_1\oplus G_0 \xrightarrow{\text{proj}_1} A_1\). We have an obvious map of complexes \(\tilde \chi: A^\chi_\bullet \to G_\bullet\) whose mapping cone \(0\to A_0 \to A_1\oplus G_0 \to G_1 \to 0\) is exact because \(\chi\) induces a surjection on the first cohomology of \(G_\bullet\), hence \(A^\chi_\bullet\) is quasi-isomorphic to \(G_\bullet\). The complex \(G_\bullet\) is perfect by definition --- it is quasi-isomorphic to \(E_\bullet\), a bounded complex of locally free sheaves --- so \(A^\chi_\bullet\) is perfect as well. Whenever the second term of a perfect two-term complex is locally free, so is its first term, and therefore \(A^\chi_\bullet\) consists of locally free sheaves as desired.

  The sheaf \(E_1\oplus F_1\) admits two obvious maps to \(G_1\) inducing surjections onto the first cohomology of \(G_1\). They are \(\phi\oplus 0\) and \(0\oplus \psi\), and using the process above, we obtain identities \(e(E_\bullet) = e(A^{\phi\oplus 0}_\bullet)\) and \(e(F_\bullet) = e(A^{0\oplus\psi}_\bullet)\) respectively. The locus of maps in the vector space \(\Hom(E_1\oplus F_1, G_1)\) inducing a surjection in cohomology is open in the Zariski topology, therefore, we can find a family of Chow classes parameterized by some \(U\subset \Hom(E_1\oplus F_1, G_1)\) containing \(e(E_\bullet)\) and \(e(F_\bullet)\). This gives us our result.
\end{proof}
